\section{Research Taste 与《金刚经》：怎样判断问题值得做}

访谈中间长段讨论 research taste，并提到《金刚经》里的“如梦幻泡影”。这部分不是玄学插曲，而是研究方法论。谢赛宁关心的是：在高度不确定、评价噪声巨大、短期结果经常误导的研究环境里，怎样保持对问题本身的判断。

\begin{figure}[H]
\centering
\includegraphics[width=0.92\textwidth]{research-taste-map.png}
\caption{Research taste 的构成：问题选择、审美、证据和时机。}
\end{figure}

\begin{knowledgebox}{读图：research taste 不是灵感，而是长期判断系统}
图中四个角分别是问题选择、审美、证据和时机。一个研究者要判断什么问题值得做，不能只看当前热度，也不能只看能否快速发论文；还要看这个问题十年后是否仍有意义、证据是否能支持方向、组织资源是否足够、自己是否能忍受长期不确定性。
\end{knowledgebox}

\subsection{审美：反对廉价引用和口号}

上一段把 research taste 拆成问题、证据、审美和时机；这里先展开“审美”。它不是文风偏好，而是研究者对概念滥用、口号包装和证据偷换的敏感度。

结尾处谢赛宁吐槽把维特根斯坦“语言的边界就是世界的边界”直接拿来给 LLM 背书，也吐槽把费曼的 “What I cannot create, I do not understand” 粗暴拿来给 unified model 背书。他不是反对哲学或名言，而是反对脱离上下文的装饰性引用。

这背后是 research taste 的审美标准：一个论点不能靠名人名言支撑，而要靠定义、假设、机制和证据。哲学可以启发研究，但不能替代 technical argument。

\begin{warningbox}{名人名言不是技术论证}
把一句哲学或物理学名言贴在论文开头，不能自动证明模型路线正确。高质量研究需要说明：概念在原语境中是什么意思，迁移到 AI 后假设是否仍成立，哪些机制和证据支持这个迁移。
\end{warningbox}

\subsection{长期主义：从拒稿到时间检验}

访谈中还讲到一篇论文曾因细节被拒，后来在另一个会议发表并获得 test-of-time award。这个故事说明，研究价值和评审结果并不总是同步。短期评审可能过度关注格式、细节或当时的主流偏好；长期影响更取决于问题是否真实、方法是否可复用、思想是否进入后续工作。

\begin{importantbox}{研究评价的时间尺度}
短期评价看 accept / reject，长期评价看这个想法是否改变问题定义、方法工具或社区语言。真正的 research taste 是在短期噪声中坚持长期可复用的方向。
\end{importantbox}

\subsection{本章小结}

Research taste 是本期技术线和个人线的粘合剂。它解释了为什么谢赛宁会选择 vision、选择 FAIR、离开确定路径、创业做 world model，也解释了他为什么反感被口号和热词牵着走。

